\section{Stable binding and changing visibility}
\label{sec:disclosure}

A commitment fixes a value before its eventual recipients can read it.
Publication changes access to that value. We give a finite phase calculus
whose views are derived from its audience annotations, followed by a
two-execution hiding theorem. This is an ideal service, using the
fixed-payload/open pattern of Canetti and Fischlin's commitment functionality
\cite{CanettiFischlin2001}. Cryptographic realization and adversarial
scheduling require additional semantics.

\subsection{Protocols, phases, and local contexts}

A disclosure protocol has the grammar
\[
 P::=\mathsf{done}
   \mid \Bind_r(c:A;S).P
   \mid \Pub_r(c;T).P.
\]
The first constructor binds a fresh cell variable $c$ in the continuation.
The second refers to an existing cell; it binds no new variable.
The finite command order, cell names, owners, audience sets, and declaration
syntax are fixed by the protocol. Data values do not control this schedule.
For each bind, $r\in S\subseteq\Roles$; for each publication, only the
recorded owner may authorize it.

A phase context $\Xi$ is an ordered list of declarations
$c:A@r,S_c$. Its underlying data context contains every declared cell.
Let $\Gamma_0$ be an initial public context, held fixed throughout the
comparison of executions. Define the local context $\Gamma_i(\Xi)$ by
keeping $\Gamma_0$ and exactly the cells for which $i\in S_c$, in their
original order. A phase is well formed when every such projected context
is a well-formed context of $\D$. Thus a visible cell's type may depend
only on earlier cells visible to that recipient.
For a declared cell $c$, write $\Gamma_i^{<c}(\Xi)$ for the projection
of the declarations strictly preceding $c$, together with $\Gamma_0$.
This is a prefix of the recipient's current local context.

At a bind occurrence $m$, an implementation supplies a pure data term
$t_m:A$. More precisely, it supplies a dependency-closed subcontext
$\Gamma_{E_m}$ of $\Gamma_r(\Xi)$ with
$\Gamma_{E_m}\vdash t_m:A$ and $\Gamma_{E_m}\vdash A:\U$.
Here $E_m$ is the set of cell names in that subcontext; public parameters
remain in $\Gamma_0$. This support need not be minimal.
Using all of $\Gamma_r(\Xi)$ always gives a support when the term is
well typed. Abstracting $t_m$ over this context gives the actor's decision
function on the corresponding local data. All private inputs used by a
decision must be represented in its context.

The phase rules are as follows.
A bind extends $\Xi$ by $c:A@r,S$, provided $c$ is fresh, $r\in S$,
the owner's implementation has the type just stated, and
$\Gamma_i(\Xi)\vdash A:\U$ for every $i\in S$.
A publication changes $S_c$ to $T$, provided $S_c\subseteq T$,
$c$ is still pending, its recorded owner is the actor, and
$\Gamma_i^{<c}(\Xi)\vdash A:\U$ for every newly added
$i\in T\setminus S_c$.
Thus the type is checked at the position where $c:A$ will be inserted,
not against declarations introduced later.
At $\mathsf{done}$ no cell may remain pending.
These conditions give a formation judgment for a protocol by processing
its commands from left to right.

The publication rule matters for dependent types. Suppose Alice privately
binds $n:\Nat_{>0}$ and then $v:\mathsf{Vec}(\Bool,n)$. She may publish $n$
and then $v$ to Bob. Publishing $v$ first is rejected: its declared type is
not yet meaningful in Bob's local context. An existential package could
instead transmit the index with its vector, but that is an explicit change
of interface. Toninho and Yoshida's projection construction addresses
unknown dependencies by bundling them and using singleton types
\cite{ToninhoYoshida2017}; the phase discipline here uses explicit staged
visibility.

\subsection{Histories and operational semantics}

A store $\rho$ assigns each declared cell a closed data value, at the type
obtained by substituting the earlier stored values in its declaration.
Let $\Omega$ be the pending cells. A configuration consists of a residual
protocol, its phase context, $\rho$, and $\Omega$.
For a fixed implementation, the next bind evaluates $t_m\rho$ in $\D$
to $a$, and makes the transition
\[
 (\Xi,\rho,\Omega)
 \xrightarrow{\Bind_r(c,a;S)}
 (\Xi,c:A@r,S;\ \rho[c\mapsto a],\ \Omega\cup\{c\}).
 \tag{20}
\]
The semicolon in the target separates the extended context from its store.
Publication makes the transition
\[
 (\Xi,\rho,\Omega)
 \xrightarrow{\Pub_r(c;T)}
 (\Xi[c.S:=T],\rho,\Omega\setminus\{c\}).
 \tag{21}
\]
These transitions advance the residual protocol by one command. Their side
conditions are the phase rules above. A history at phase $m$ is a legal
sequence of the first $m$ transitions, including their intermediate stores.
The only successful terminal history has residual $\mathsf{done}$ and
$\Omega=\varnothing$.
Strong normalization and canonical forms in $\D$ make the pure evaluation
in (20) terminate. There is no other source of scheduling nondeterminism.

The audience changes in (21) are monotone. Stored values remain present
even after their pending obligations have been resolved. In particular,
the cell name used as an index in a dependent type remains the same name
across publication; (21) introduces no fresh witness.

\begin{lemma}[Phase preservation]
\label{lem:phase-preservation}
A well-formed protocol implementation and a well-typed store advance to a
well-formed phase and a well-typed store.
\end{lemma}
\begin{proof}
At a bind, substitution types $t_m\rho$ at $A\rho$.
For each recipient, $A$ is well formed in its previous local context, so
context extension by $c:A$ is valid. Other local contexts are unchanged.
At publication the data store does not change. For a new recipient write
its old local context as $\Gamma_i^{<c}(\Xi),\Delta_i$.
The publication premise types $A$ in that prefix. Structural context
insertion in $\D$ therefore gives the well-formed context
$\Gamma_i^{<c}(\Xi),c:A,\Delta_i$, preserving the judgments for the
later declarations. This is exactly the new projected context, in the
original declaration order. The underlying data context and store are
unchanged; other recipients keep their contexts. No strengthening from
the full current recipient context is required.
\end{proof}

\begin{lemma}[Binding across arbitrary continuations]
\label{prop:binding}
Fix a legal history ending with the binding of $c$ to $a$. Any two legal
extensions of that history, even using different later implementations,
publish $a$ whenever they successfully publish $c$.
\end{lemma}
\begin{proof}
Every later bind allocates a fresh name. Every publication changes only
audience metadata and pending status. Induction on each extension therefore
gives $\rho(c)=a$ at every later phase. The value exposed by a successful
publication is exactly this stored value. No preservation condition on the
later implementations is assumed.
\end{proof}

\subsection{Views derived from audiences}

At a fixed phase, define
\[
 V_i^m(h)=\bigl(m,\ (\rho_h(c))_{c\in\Gamma_i(\Xi_m)\setminus\Gamma_0}\bigr).
 \tag{22}
\]
The fixed public context and declaration metadata are implicit in this
tuple. Unavailable payloads and their instantiated private indices are
absent. The tuple is ordered by the declarations in $\Xi_m$.
Because cells are retained and audiences only expand, all previously
received values and all previous owned choices remain in the view.
This is perfect recall for this fixed-schedule language.

Thus observations are determined by the protocol and its audiences.
A decision term at phase $m$ can use only the corresponding projected
context. Histories with equal views give equal arguments to the same
implementation, hence convertible decisions in the common answer type.
The answer type is meaningful locally by the phase rules.

\begin{example}[A sealed pair]
Use distinct role names $\mathsf{Alice},\mathsf{Bob}$ and public payload
types $X,Y$. The protocol is
\[
 \Bind_{\mathsf{Alice}}(c:X;\{\mathsf{Alice}\}).
 \Bind_{\mathsf{Bob}}(d:Y;\{\mathsf{Bob}\}).
 \Pub_{\mathsf{Alice}}(c;\Roles).
 \Pub_{\mathsf{Bob}}(d;\Roles).
 \mathsf{done}.
 \tag{23}
\]
At Bob's binding phase his projected context contains no $c$.
For a fixed Bob implementation and fixed other inputs, changing Alice's
chosen payload therefore cannot change Bob's choice.
If Alice's publication is moved before Bob's bind, his context contains
$c:X$ and a policy depending on that value becomes well typed.
\end{example}

\subsection{Hiding, including dependent intermediate data}
\label{sec:hiding}

It would be false to infer secrecy until $\Pub(c)$ merely from
$i\notin S_c$. An authorized reader may bind a copy $d:=c$ and publish
$d$ first. This is a different publication of information about $c$.
The following dependency discipline states exactly which prefixes exclude
such an indirect release.

Fix an observer $i$ just after a private cell $c$ has been bound.
Assume $i\notin S_c$, and initially mark $D=\{c\}$.
Process a continuation prefix as follows:
\begin{itemize}
\item At a bind of $d$ with support $E_m$, mark $d$ precisely when
      $E_m\cap D\ne\varnothing$, retaining all previous marks.
      If $d$ is marked, require $i\notin S_d$.
\item At a publication of $d$ to $T$, require $i\notin T$ whenever $d\in D$.
\end{itemize}
Call such a prefix \emph{sealed for $i$ from $c$}. These are decidable
side conditions on the supplied supports and the fixed audience annotations.
A derivation in a smaller support can certify a prefix that a coarser
dependency analysis rejects. Types are included in the support judgment,
so marking tracks dependencies of answer types as well as of term bodies.

\begin{theorem}[Two-execution hiding for sealed prefixes]
\label{thm:hiding}
Start from the same history before the binding of $c$, and bind it to any
two values $a,a':A$ with the same initial audience excluding $i$.
Run the same well-typed continuation implementations along a prefix sealed
for $i$ from $c$.
At every phase of that prefix, all unmarked cells have convertible values
in the same instantiated types, and the two $i$-views agree.
Consequently every decision made by $i$ in the prefix is independent of
whether $a$ or $a'$ was stored.
\end{theorem}
\begin{proof}
Immediately after the two binds, all old cells agree, and the only marked
cell is invisible to $i$. Maintain these two invariants by induction.

For a bind with support disjoint from $D$, the two substitutions agree on
its whole typing context. Congruence and substitution in $\D$ give the same
answer type and convertible evaluations of $t_m$. Its newly stored value is
therefore unmarked and equal in both runs. If the support meets $D$, mark
the new cell. Its values and even its instantiated types may differ, but
the sealing condition excludes it from $i$'s context. It cannot invalidate
the equality claim for unmarked cells.

A publication changes no value. If it exposes a cell to $i$, that cell is
unmarked by the sealing condition and its values already agree. Otherwise
it changes neither $i$'s data nor the equality invariant. All phase metadata
are the same in the two runs, since the schedule and audience annotations
do not depend on payloads. Equation (22) now gives equal views.
The observer's next well-typed decision is an evaluation in that same
visible context and so also agrees.
\end{proof}

This is a noninterference statement with an explicit release boundary.
It proves hiding until the first release to $i$ of $c$ or any cell
depending on it. Hiding until the particular command $\Pub(c)$ follows
when no earlier dependent release occurs. The copy-and-publish example
shows why that qualification is necessary. For (23), Bob's bind has a
support disjoint from $\{c\}$; the theorem supplies the sealed-pair claim
from the audience rules.

The result uses ordinary dependent substitution and dependency analysis.
Its role in the proposed calculus is to specify an obligation for
information-aware connection: a context must preserve the sealing judgment,
rather than only the payload types.

\subsection{Composition and common knowledge}

For a map $u$ from one component's histories to another's, define
observation-respecting connection by the condition
\[
 V_i(h)=V_i(h')\ \Longrightarrow\
 V'_i(u(h))=V'_i(u(h')).
 \tag{24}
\]
Here both view maps are computed from their protocols. A typed connection
also transports the corresponding dependent answer families.

\begin{lemma}[Composition of view-respecting maps]
\label{lem:view-composition}
Such maps include identities, are closed under composition, and preserve
the uniformity of a policy pulled back along them.
\end{lemma}
\begin{proof}
Apply (24) successively to the maps. An equal target view supplies the same
arguments to a fixed target policy.
\end{proof}

At a fixed public stage let $W$ be the possible histories, with
indistinguishability given by equality of $V_i$. Common knowledge among a
group $S$ means truth at every history reachable by a finite chain of these
relations. This is the standard partition characterization
\cite{Aumann1976,HalpernMoses1990}.

\begin{lemma}[Characterization of common readout knowledge]
\label{thm:common-knowledge}
For $p:W\to P$, the fact $p=p(h)$ is common knowledge among $S$ at every
$h\in W$ if and only if
\[
 \forall i\in S,\quad
 V_i(h)=V_i(h')\ \Longrightarrow\ p(h)=p(h').
 \tag{25}
\]
\end{lemma}
\begin{proof}
For sufficiency, each indistinguishability edge preserves $p$, hence so
does every finite chain. For necessity, any one such edge is a chain;
common knowledge at its source gives the required equality at its target.
\end{proof}

After successful publication to $S$, the cell name and value form such a
readout. The public-stage assumption includes common availability of the
stage itself; independent delayed message delivery requires a different
model of views \cite{HalpernMoses1990,BaltagMossSolecki1999}.

\subsection{Failure and endpoint direction}

With a coproduct, a fallible resolution can have interface
\[
 \Unit+\sum_{a:A}\Opens(c,a).
 \tag{26}
\]
The right injection supplies the opened payload and evidence that it is the
value fixed in $c$; the left reports failure. Only success adds the payload
to the recipient's data context. The successful, fixed-schedule semantics
above does not give a progress or hiding theorem for data-dependent failure,
withholding, or timing. Those events can themselves disclose information and
must appear in the views of an extended calculus.

VegasCore stores immutable private commitments, exposes success or failure,
and indexes programs by outstanding commitments
\cite{VegasCore2026}. Its implementation motivates this distinction between
successful disclosure and resolution of an obligation.

Publication has a producing endpoint and receiving endpoints, giving the
usual input/output polarity \cite{Griffith2016}. Binding and publication
are temporal phases of the same stable cell. They need not exchange owners,
and their composition is not an involution.
